
sur @@M0@@, car ces
points sont en correspondance biunivoque avec les drapeaux du type @@M1@@ (@@M2@@ termes) de l'espace vectoriel @@M3@@ de dimension @@M4@@ sur @@M5@@. D'autre part ces points rationnels correspondent aussi aux germes de sections de @@M6@@ au voisinage de @@M7@@, de sorte qu'on peut trouver une telle section @@M8@@ au-dessus d'un ouvert @@M9@@ de @@M10@@. Soit @@M11@@ l'adhérence de @@M12@@ munie de la structure induite réduite, et soit @@M13@@ la restriction de @@M14@@ à @@M15@@. Le morphisme @@M16@@ est projectif, et comme il est un isomorphisme sur @@M17@@, il est birationnel.
Donc @@M18@@ (EGA IV 5.6.6.1). D'autre part @@M19@@ est fermé et contient @@M20@@, de sorte que @@M21@@ est surjectif et @@M22@@ (EGA IV 5.4.1(ii)).

Les hypothèses du cas I sont donc satisfaites pour @@M23@@ et @@M24@@, et on trouve
@@M25@@
On en conclut, tenant compte de la formule de projection et de (1.1.8) :
@@M26@@

D'autre part l'homomorphisme canonique
@@M27@@
est un isomorphisme sur l'ouvert dense @@M28@@, et les @@M29@@-modules @@M30@@ (@@M31@@) sont concentrés sur le sous-ensemble fermé rare @@M32@@. Il s'ensuit que
\begin{equation}\tag{1.3.9}
f_*(\cl_\bullet(\underline{O}_{X'}))=\cl_\bullet(\underline{O}_X)+w,
